\chapter{2024年全国甲卷(理)}

\section{单选题}

\begin{problem}
设 \( z=5+\i \),则 \( \i \left( \bar z+z \right)= \)(\quad)

\choices
    {\( 10\i \)}
    {\( 2\i \)}
    {\( 10 \)}
    {\( 2 \)}
\end{problem}

\begin{problem}
集合 \( A=\{1,2,3,4,5,9\} \),\( B=\left\{ x \mid \sqrt{x}\in A \right\} \),则 \( C_A\left( A\cap B \right)= \)(\quad)

\choices
    {\( \{1,4,9\} \)}
    {\( \{3,4,9\} \)}
    {\( \{1,2,3\} \)}
    {\( \{2,3,5\} \)}
\end{problem}

\begin{problem}
若实数 \(x\),\(y\) 满足约束条件
\(\begin{cases}
4x-3y-3\ge 0,\\
x-2y-2\le 0,\\
2x+6y-9\le 0,
\end{cases}\)
则 \( z=x-5y \) 的最小值为(\quad)

\choices
    {\( \dfrac{1}{2} \)}
    {\( 0 \)}
    {\( -\dfrac{5}{2} \)}
    {\( -\dfrac{7}{2} \)}
\end{problem}

\begin{problem}
等差数列 \( \{a_n\} \) 的前 \( n \) 项和为 \( S_n \),若 \( S_5=S_{10} \),\( a_5=1 \),则 \( a_1= \)(\quad)

\choices
    {\( \dfrac{7}{2} \)}
    {\( \dfrac{7}{3} \)}
    {\( -\dfrac{1}{3} \)}
    {\( -\dfrac{7}{11} \)}
\end{problem}

\begin{problem}
已知双曲线的两个焦点分别为 \( (0,4) \),\( (0,-4) \),点 \( P(-6,4) \) 在该双曲线上,则该双曲线的离心率为(\quad)

\choices
    {\( 4 \)}
    {\( 3 \)}
    {\( 2 \)}
    {\( \sqrt{2} \)}
\end{problem}

\begin{problem}
设函数 \( f(x)=\dfrac{\e^x+2\sin x}{1+x^2} \),则曲线 \( y=f(x) \) 在 \( (0,1) \) 处的切线与两坐标轴围成的三角形的面积为(\quad)

\choices
    {\( \dfrac{1}{6} \)}
    {\( \dfrac{1}{3} \)}
    {\( \dfrac{1}{2} \)}
    {\( \dfrac{2}{3} \)}
\end{problem}

\begin{problem}
函数 \( f(x)=-x^2+\left( \e^x-\e^{-x} \right)\sin x \) 在区间 \( [-2.8,2.8] \) 的大致图像为(\quad)

\choices
    {\choicebitmap[width=0.15\paperwidth]{img/2024/national_paper_a_science/q7_optA.png}}
    {\choicebitmap[width=0.15\paperwidth]{img/2024/national_paper_a_science/q7_optB.png}}
    {\choicebitmap[width=0.15\paperwidth]{img/2024/national_paper_a_science/q7_optC.png}}
    {\choicebitmap[width=0.15\paperwidth]{img/2024/national_paper_a_science/q7_optD.png}}
\end{problem}

\begin{problem}
已知 \( \dfrac{\cos \alpha}{\cos \alpha-\sin \alpha}=\sqrt{3} \),则 \( \tan \left( \alpha+\dfrac{\pi}{4} \right)= \)(\quad)

\choices
    {\( 2\sqrt{3}+1 \)}
    {\( 2\sqrt{3}-1 \)}
    {\( \dfrac{\sqrt{3}}{2} \)}
    {\( 1-\sqrt{3} \)}
\end{problem}

\begin{problem}
已知向量 \( \bs a=(x+1,x) \),\( \bs b=(x,2) \),则(\quad)

\choices
    {``\( x=-3 \)''是``\( \bs a\perp \bs b \)''的必要条件}
    {``\( x=1+\sqrt{3} \)''是``\( \bs a\parallel\bs b \)''的必要条件}
    {``\( x=0 \)''是``\( \bs a\perp \bs b \)''的充分条件}
    {``\( x=-1+\sqrt{3} \)''是``\( \bs a\parallel\bs b \)''的充分条件}
\end{problem}

\begin{problem}
设 \(\alpha\),\(\beta\) 是两个平面,\(m\),\(n\) 是两条直线,且 \( \alpha\cap\beta=m \). 下列四个命题:

\begin{circlelist}
\item 若 \( m\parallel n \),则 \( n\parallel\alpha \) 或 \( n\parallel\beta \);
\item 若 \( m\perp n \),则 \( n\perp \alpha \) 或 \( n\perp \beta \);
\item 若 \( n\parallel\alpha \),且 \( n\parallel\beta \),则 \( m\parallel n \);
\item 若 \( n \) 与 \( \alpha \) 和 \( \beta \) 所成的角相等,则 \( m\perp n \).
\end{circlelist}

其中所有真命题的编号是(\quad)

\choices
    {\( \circled{1}\circled{3} \)}
    {\( \circled{2}\circled{4} \)}
    {\( \circled{1}\circled{2}\circled{3} \)}
    {\( \circled{1}\circled{3}\circled{4} \)}
\end{problem}

\begin{problem}
在 \( \triangle ABC \) 中内角 \(A\),\(B\),\(C\) 所对边分别为 \(a\),\(b\),\(c\),若 \( B=\dfrac{\pi}{3} \),\( b^2=\dfrac{9}{4}ac \),则 \( \sin A+\sin C= \)(\quad)

\choices
    {\( \dfrac{3}{2} \)}
    {\( \sqrt{2} \)}
    {\( \dfrac{\sqrt{7}}{2} \)}
    {\( \dfrac{\sqrt{3}}{2} \)}
\end{problem}

\begin{problem}
已知 \( b \) 是 \(a\),\(c\) 的等差中项,直线 \( ax+by+c=0 \) 与圆 \( x^2+y^2+4y-1=0 \) 交于 \(A\),\(B\) 两点,则 \( \left| AB \right| \) 的最小值为(\quad)

\choices
    {\( 1 \)}
    {\( 2 \)}
    {\( 4 \)}
    {\( 2\sqrt{5} \)}
\end{problem}

\section{填空题}

\begin{problem}
\( \left( \dfrac{1}{3}+x \right)^{10} \) 的展开式中,各项系数中的最大值是\fillinblank{}.
\end{problem}

\begin{problem}
已知圆台甲,乙的上底面半径均为 \( r_1 \),下底面半径均为 \( r_2 \),母线长分别为 \( 2\left( r_2-r_1 \right) \) 和 \( 3\left( r_2-r_1 \right) \),则圆台甲与乙的体积之比 \( \dfrac{V_{\text{甲}}}{V_{\text{乙}}}= \)\fillinblank{}.
\end{problem}

\begin{problem}
已知 \( a>1 \),\( \dfrac{1}{\log_8 a}-\dfrac{1}{\log_a 4}=-\dfrac{5}{2} \),则 \( a= \)\fillinblank{}.
\end{problem}

\begin{problem}
有 \( 6\text{ 个} \) 相同的球,分别标有数字 \(1\),\(2\),\(3\),\(4\),\(5\),\(6\),从中不放回地随机抽取 \( 3\text{ 次} \),每次取 \( 1\text{ 个} \) 球. 记 \( m \) 为前两次取出的球上数字的平均值,\( n \) 为取出的三个球上数字的平均值,则 \( m \) 与 \( n \) 差的绝对值不超过 \( \dfrac{1}{2} \) 的概率是\fillinblank{}.
\end{problem}

\section{解答题}

\begin{problem}
某工厂进行生产线智能化升级改造,升级改造后,从该工厂甲,乙两个车间的产品中随机抽取 \( 150\text{ 件} \) 进行检验,数据如下:

\begin{center}
\renewcommand{\arraystretch}{1.3}
\begin{tabular}{c|c|c|c|c}
 & 优级品 & 合格品 & 不合格品 & 总计 \\
\hline
甲车间 & 26 & 24 & 0 & 50 \\
\hline
乙车间 & 70 & 28 & 2 & 100 \\
\hline
总计 & 96 & 52 & 2 & 150
\end{tabular}
\end{center}

\begin{enumerate}
\item 填写如下列联表:

\begin{center}
\renewcommand{\arraystretch}{1.3}
\begin{tabular}{c|c|c}
 & 优级品 & 非优级品 \\
\hline
甲车间 &\fillinblank{}&\fillinblank{} \\
\hline
乙车间 &\fillinblank{}&\fillinblank{}
\end{tabular}
\end{center}

能否有 \( 95\% \) 的把握认为甲,乙两车间产品的优级品率存在差异?能否有 \( 99\% \) 的把握认为甲,乙两车间产品的优级品率存在差异?
\item 已知升级改造前该工厂产品的优级品率 \( p=0.5 \),设 \( \bar p \) 为升级改造后抽取的 \( n\text{ 件} \) 产品的优级品率. 如果 \( \bar p>p+1.65\sqrt{\frac{p(1-p)}{n}} \),则认为该工厂产品的优级品率提高了. 根据抽取的 \( 150\text{ 件} \) 产品的数据,能否认为生产线智能化升级改造后,该工厂产品的优级品率提高了?(\( \sqrt{150}\approx 12.247 \))
\end{enumerate}

附:\( K^2=\frac{n(ad-bc)^2}{(a+b)(c+d)(a+c)(b+d)} \)

\begin{center}
\renewcommand{\arraystretch}{1.3}
\begin{tabular}{c|c|c|c}
\( P\left( K^2\ge k \right) \) & 0.050 & 0.010 & 0.001 \\
\hline
\( k \) & 3.841 & 6.635 & 10.828
\end{tabular}
\end{center}
\end{problem}

\begin{problem}
记 \( S_n \) 为数列 \( \{a_n\} \) 的前 \( n \) 项和,且 \( 4S_n=3a_n+4 \).

\begin{enumerate}
\item 求 \( \{a_n\} \) 的通项公式;
\item 设 \( b_n=\left( -1 \right)^{n-1}na_n \),求数列 \( \{b_n\} \) 的前 \( n \) 项和 \( T_n \).
\end{enumerate}
\end{problem}

\begin{problem}
如图,在以 \(A\),\(B\),\(C\),\(D\),\(E\),\(F\) 为顶点的五面体中,四边形 \(ABCD\) 与四边形 \(ADEF\) 均为等腰梯形,\( BC\parallel AD \),\( EF\parallel AD \),\( AD=4 \),\( AB=BC=EF=2 \),\( ED=\sqrt{10} \),\( FB=2\sqrt{3} \),\(M\) 为 \(AD\) 的中点.

\begin{enumerate}
\item 证明:\( BM\parallel \)平面 \(CDE\);
\item 求二面角 \( F-BM-E \) 的正弦值.
\end{enumerate}

\bitmapfigure[max width=0.55\linewidth]{img/2024/national_paper_a_science/q19_fig1.png}
\end{problem}

\begin{problem}
设椭圆 \( C:\dfrac{x^2}{a^2}+\dfrac{y^2}{b^2}=1 \)(\( a>b>0 \))的右焦点为 \(F\),点 \( M\left( 1,\dfrac{3}{2} \right) \) 在 \(C\) 上,且 \( MF\perp x \) 轴.

\begin{enumerate}
\item 求 \(C\) 的方程;
\item 过点 \( P(4,0) \) 的直线与 \(C\) 交于 \(A\),\(B\) 两点,\(N\) 为线段 \(FP\) 的中点,直线 \(NB\) 交直线 \(MF\) 于点 \(Q\),证明:\( AQ\perp y \) 轴.
\end{enumerate}
\end{problem}

\begin{problem}
已知函数 \( f(x)=(1-ax)\ln(1+x)-x \).

\begin{enumerate}
\item 当 \( a=-2 \) 时,求 \( f(x) \) 的极值;
\item 当 \( x\ge 0 \) 时,\( f(x)\ge 0 \) 恒成立,求 \( a \) 的取值范围.
\end{enumerate}
\end{problem}

\section{选考题:解答题}

\begin{problem}
在平面直角坐标系 \( xOy \) 中,以坐标原点 \(O\) 为极点,\( x \) 轴的正半轴为极轴建立极坐标系,曲线 \(C\) 的极坐标方程为 \( \rho=\rho\cos\theta+1 \).

\begin{enumerate}
\item 写出 \(C\) 的直角坐标方程;
\item 设直线 \( l \) 的参数方程为 \(\begin{cases}
x=t,\\
y=t+a,
\end{cases}\)(\( t \) 为参数),若 \(C\) 与 \( l \) 相交于 \(A\),\(B\) 两点,且 \( |AB|=2 \),求 \( a \) 的值.
\end{enumerate}
\end{problem}

\begin{problem}
实数 \(a\),\(b\) 满足 \( a+b\ge 3 \).

\begin{enumerate}
\item 证明:\( 2a^2+2b^2>a+b \);
\item 证明:\( |a-2b^2|+|b-2a^2|\ge 6 \).
\end{enumerate}
\end{problem}

